\documentclass[12pt, a4paper]{article}

% encoding y lenguaje
\usepackage[utf8]{inputenc}
\usepackage[T1]{fontenc}
\usepackage[spanish]{babel}

% layout
\usepackage[top=2.5cm, bottom=2.5cm, left=3cm, right=3cm]{geometry}
\usepackage{parskip}
\usepackage{microtype}

% tipografía
\usepackage{lmodern}

% código
\usepackage{listings}
\usepackage{xcolor}

% tablas
\usepackage{booktabs}
\usepackage{tabularx}
\usepackage{array}
\usepackage{multirow}

% gráficos / diagramas
\usepackage{tikz}
\usetikzlibrary{shapes.geometric, arrows.meta, positioning, fit, backgrounds}

% misc
\usepackage{hyperref}
\usepackage{enumitem}
\usepackage{fancyhdr}
\usepackage{titlesec}

% ──────────────────────────────────────────────
% colores del design system de Caché
% ──────────────────────────────────────────────
\definecolor{ink}{HTML}{0A0A0A}
\definecolor{acid}{HTML}{D4FF1A}
\definecolor{blood}{HTML}{FF2E2E}
\definecolor{pulse}{HTML}{00FF88}
\definecolor{bone}{HTML}{E8E6DF}
\definecolor{softgray}{HTML}{6B6B6B}
\definecolor{codebg}{HTML}{1E1E1E}
\definecolor{codestr}{HTML}{CE9178}
\definecolor{codecomment}{HTML}{6A9955}
\definecolor{codekw}{HTML}{569CD6}
\definecolor{codeann}{HTML}{DCDCAA}
\definecolor{neo4jcolor}{HTML}{018BFF}
\definecolor{mongocolor}{HTML}{00ED64}
\definecolor{rediscolor}{HTML}{FF4438}
\definecolor{casscolor}{HTML}{1287B1}

% ──────────────────────────────────────────────
% estilo de código
% ──────────────────────────────────────────────
\lstdefinestyle{java}{
  backgroundcolor=\color{codebg},
  basicstyle=\ttfamily\footnotesize\color{white},
  keywordstyle=\color{codekw}\bfseries,
  stringstyle=\color{codestr},
  commentstyle=\color{codecomment}\itshape,
  numberstyle=\tiny\color{softgray},
  breaklines=true,
  breakatwhitespace=true,
  frame=single,
  framerule=0.5pt,
  rulecolor=\color{softgray},
  numbers=left,
  numbersep=8pt,
  tabsize=4,
  showstringspaces=false,
  language=Java,
  morekeywords={String, List, Instant, BigDecimal, GeoJsonPoint, Optional},
  xleftmargin=1.5em,
  xrightmargin=0.5em,
}

\lstdefinestyle{cypher}{
  backgroundcolor=\color{codebg},
  basicstyle=\ttfamily\footnotesize\color{white},
  keywordstyle=\color{codekw}\bfseries,
  stringstyle=\color{codestr},
  commentstyle=\color{codecomment}\itshape,
  breaklines=true,
  frame=single,
  framerule=0.5pt,
  rulecolor=\color{softgray},
  tabsize=2,
  showstringspaces=false,
  morekeywords={MATCH, RETURN, WHERE, ORDER, BY, LIMIT, CREATE, WITH, AS, DESC, ASC, EXISTS, NOT, AND, OR, IN, OPTIONAL},
  xleftmargin=1.5em,
}

\lstdefinestyle{redis}{
  backgroundcolor=\color{codebg},
  basicstyle=\ttfamily\footnotesize\color{white},
  keywordstyle=\color{rediscolor}\bfseries,
  commentstyle=\color{codecomment}\itshape,
  breaklines=true,
  frame=single,
  framerule=0.5pt,
  rulecolor=\color{softgray},
  xleftmargin=1.5em,
}

% ──────────────────────────────────────────────
% encabezado / pie de página
% ──────────────────────────────────────────────
\pagestyle{fancy}
\fancyhf{}
\fancyhead[L]{\small\texttt{caché — arquitectura de datos}}
\fancyhead[R]{\small\texttt{v0.1 / 2025}}
\fancyfoot[C]{\small\thepage}
\renewcommand{\headrulewidth}{0.4pt}

% ──────────────────────────────────────────────
% títulos
% ──────────────────────────────────────────────
\titleformat{\section}{\large\bfseries}{}{0em}{}[\titlerule]
\titleformat{\subsection}{\normalsize\bfseries}{}{0em}{}

% ──────────────────────────────────────────────
% comandos utilitarios
% ──────────────────────────────────────────────
\newcommand{\dbname}[2]{\textbf{\textcolor{#1}{#2}}}
\newcommand{\code}[1]{\texttt{#1}}
\newcommand{\pk}[1]{\underline{\texttt{#1}}}   % primary key
\newcommand{\ck}[1]{\texttt{\textit{#1}}}       % clustering key
\newcommand{\idx}[1]{\texttt{#1}\textsuperscript{*}} % indexed

% ──────────────────────────────────────────────

\begin{document}

% ──────────────────────────────────────────────
% PORTADA
% ──────────────────────────────────────────────
\begin{titlepage}
  \centering
  \vspace*{3cm}
  {\Huge\bfseries caché}\\[0.5em]
  {\large arquitectura de bases de datos}\\[2em]
  \rule{0.6\textwidth}{0.5pt}\\[2em]
  {\normalsize
    Red social de descubrimiento nocturno.\\
    Arquitectura políglota: Neo4j · MongoDB · Redis · Cassandra.\\[1em]
    Cada motor fue elegido por lo que hace mejor,\\
    no por uniformidad tecnológica.
  }\\[3em]
  \rule{0.6\textwidth}{0.5pt}\\[2em]
  {\small\texttt{backend: Spring Boot 3.3 / Java 21}}\\
  {\small\texttt{frontend: Next.js 14}}\\[2em]
  {\footnotesize Buenos Aires, 2025}
\end{titlepage}

\tableofcontents
\newpage

% ──────────────────────────────────────────────
\section{Contexto y estrategia de datos}
% ──────────────────────────────────────────────

Caché es una app donde los usuarios descubren y se anotan en eventos nocturnos.
Las operaciones centrales son:

\begin{enumerate}[nosep]
  \item Mostrar un feed personalizado de eventos (\emph{¿qué hay esta noche?})
  \item Anotar a un usuario en un evento (\emph{check-in})
  \item Mostrar presencia en tiempo real (\emph{¿quién está ahora mismo?})
  \item Recomendar eventos por red social (\emph{¿adónde va tu gente?})
  \item Analizar tendencias históricas (\emph{¿qué días y horarios movieron más?})
\end{enumerate}

Ningún motor cubre los cinco casos con la misma eficiencia.
La decisión fue asignar a cada motor el caso para el que fue diseñado:

\vspace{1em}
\begin{tabularx}{\textwidth}{@{}lXl@{}}
  \toprule
  \textbf{Motor} & \textbf{Responsabilidad} & \textbf{Fortaleza explotada} \\
  \midrule
  \dbname{neo4jcolor}{Neo4j}     & Recomendaciones por red social & Traversal de grafos \\
  \dbname{mongocolor}{MongoDB}   & Catálogo de eventos            & Schema flexible + geo \\
  \dbname{rediscolor}{Redis}     & Presencia en tiempo real       & Estructuras en memoria \\
  \dbname{casscolor}{Cassandra}  & Historial y tendencias         & Writes masivos + series temporales \\
  \bottomrule
\end{tabularx}

\vspace{1em}
El servicio \code{CheckinService} orquesta los cuatro motores en una sola operación
de negocio. La sección \ref{sec:checkin} detalla ese flujo.

% ──────────────────────────────────────────────
\section{Neo4j — grafo de relaciones}
\label{sec:neo4j}
% ──────────────────────────────────────────────

\subsection{Por qué un grafo}

Las recomendaciones de Caché son inherentemente sociales: \emph{``tu amigo Mati
se anotó en SUB00 — ¿vas?''} Ese tipo de pregunta requiere recorrer relaciones
entre entidades, no filtrar filas. Con SQL o MongoDB resolvería con JOINs
recursivos o \$lookup con múltiples niveles; con un grafo, es un pattern de
2 saltos:

\begin{lstlisting}[style=cypher, caption={consulta principal de recomendación}]
MATCH (me:User {userId: $userId})
      -[:FRIENDS_WITH]-(friend:User)
      -[:ATTENDING]->(e:Event)
WHERE e.startsAtEpoch > $nowEpoch
RETURN e.eventId AS eventId, count(friend) AS friendCount
ORDER BY friendCount DESC
LIMIT $limit
\end{lstlisting}

La complejidad de esta query es proporcional a los amigos del usuario,
no al total de usuarios en la plataforma. Eso no se puede lograr con
índices relacionales sin un grafo.

\subsection{Nodos}

\subsubsection{\code{User}}

Representa a un usuario de la app.

\begin{tabularx}{\textwidth}{@{}llX@{}}
  \toprule
  \textbf{Campo} & \textbf{Tipo} & \textbf{Descripción} \\
  \midrule
  \code{id}       & \code{Long}   & ID interno de Neo4j (generado) \\
  \code{userId}   & \code{String} & ID de la app (UUID del sistema de usuarios) \\
  \code{name}     & \code{String} & Nombre para display \\
  \code{city}     & \code{String} & Ciudad base (filtro geográfico de recomendaciones) \\
  \bottomrule
\end{tabularx}

\vspace{0.5em}
El \code{userId} es el puente con los otros motores. Neo4j solo guarda lo
mínimo necesario para el traversal; el perfil completo vive en otro sistema.

\subsubsection{\code{Event}}

Espejo liviano de los eventos de MongoDB, solo con los campos necesarios
para el grafo.

\begin{tabularx}{\textwidth}{@{}llX@{}}
  \toprule
  \textbf{Campo}        & \textbf{Tipo}   & \textbf{Descripción} \\
  \midrule
  \code{eventId}        & \code{String}   & ID del documento en MongoDB (\emph{referencia}) \\
  \code{name}           & \code{String}   & Nombre del evento \\
  \code{venueId}        & \code{String}   & ID del venue \\
  \code{genre}          & \code{String}   & Género musical (para recomendaciones por afinidad) \\
  \code{city}           & \code{String}   & Ciudad \\
  \code{startsAtEpoch}  & \code{long}     & Timestamp epoch en milisegundos \\
  \bottomrule
\end{tabularx}

\vspace{0.5em}
\textbf{¿Por qué \code{startsAtEpoch} como \code{long} y no como fecha?}
Los tipos de fecha varían entre drivers y versiones de Neo4j. Un
\code{long} es un tipo primitivo sin ambigüedad; la comparación en Cypher
es aritmética simple (\code{> \$nowEpoch}).

\subsection{Relaciones}

\subsubsection{\code{FRIENDS\_WITH}}

Relación bidireccional sin propiedades entre dos nodos \code{User}.
No tiene dirección efectiva: si A es amigo de B, se puede recorrer en ambas
direcciones con \code{-[:FRIENDS\_WITH]-} (sin flecha).

\subsubsection{\code{ATTENDING} (con propiedades)}

Relación de un \code{User} hacia un \code{Event}. Tiene propiedades propias
implementadas con \code{@RelationshipProperties}:

\begin{tabularx}{\textwidth}{@{}llX@{}}
  \toprule
  \textbf{Campo}       & \textbf{Tipo}     & \textbf{Descripción} \\
  \midrule
  \code{registeredAt}  & \code{Instant}    & Cuándo se anotó \\
  \code{status}        & \code{String}     & \code{confirmed} / \code{pending} / \code{cancelled} \\
  \bottomrule
\end{tabularx}

\vspace{0.5em}
Las propiedades en la relación evitan crear un nodo intermedio.
Un nodo \code{Attendance} sería un anti-patrón en Neo4j cuando el
único propósito es almacenar metadata de la arista.

\subsection{Diagrama del grafo}

\begin{center}
\begin{tikzpicture}[
  node distance=2.5cm,
  usernode/.style={
    ellipse, draw=neo4jcolor, fill=neo4jcolor!15,
    minimum width=2.2cm, minimum height=0.9cm, font=\small\bfseries
  },
  eventnode/.style={
    ellipse, draw=mongocolor, fill=mongocolor!15,
    minimum width=2.2cm, minimum height=0.9cm, font=\small\bfseries
  },
  rel/.style={-Stealth, thick},
  birel/.style={Stealth-Stealth, thick},
  relabel/.style={font=\scriptsize\ttfamily, midway, fill=white, inner sep=2pt}
]
  \node[usernode] (me)     {User A};
  \node[usernode] (fr1)    [right=3cm of me]  {User B};
  \node[usernode] (fr2)    [below=1.5cm of fr1] {User C};
  \node[eventnode](ev1)    [right=3cm of fr1]  {Event 1};
  \node[eventnode](ev2)    [right=3cm of fr2]  {Event 2};

  \draw[birel] (me)  -- node[relabel, above] {FRIENDS\_WITH} (fr1);
  \draw[birel] (me)  -- node[relabel, left]  {FRIENDS\_WITH} (fr2);
  \draw[rel]   (fr1) -- node[relabel, above] {ATTENDING}     (ev1);
  \draw[rel]   (fr1) -- node[relabel, below left] {ATTENDING} (ev2);
  \draw[rel]   (fr2) -- node[relabel, below] {ATTENDING}     (ev2);
\end{tikzpicture}
\end{center}

La query de recomendación parte de \emph{User A}, recorre
\code{FRIENDS\_WITH} hacia B y C, luego \code{ATTENDING} hacia los eventos.
\emph{Event 2} aparece dos veces (B y C van) → \code{friendCount = 2},
sube en el ranking.

% ──────────────────────────────────────────────
\section{MongoDB — catálogo de eventos}
\label{sec:mongo}
% ──────────────────────────────────────────────

\subsection{Por qué documentos}

Un evento nocturno tiene estructura variable: algunos tienen lineup, otros no;
algunos tienen precio libre, otros entrada fija; algunos tienen foto, otros son
flyers tipográficos. Un esquema relacional fijo requeriría nullable columns o
tablas auxiliares para cada variante. MongoDB permite que cada documento tenga
la forma que necesita.

Además, el caso de uso de búsqueda geográfica (\emph{``eventos cerca de mí''})
tiene soporte nativo con índices \code{2dsphere}, que en bases relacionales
requiere extensiones (PostGIS) o soluciones ad-hoc.

\subsection{Documento: \code{EventDocument}}

\begin{tabularx}{\textwidth}{@{}llX@{}}
  \toprule
  \textbf{Campo}       & \textbf{Tipo}            & \textbf{Descripción} \\
  \midrule
  \code{\_id}          & \code{String}            & ID ObjectId de MongoDB \\
  \idx{name}           & \code{String}            & Nombre del evento \\
  \code{venueId}       & \code{String}            & ID del venue (referencia externa) \\
  \code{venueName}     & \code{String}            & Nombre del venue (desnormalizado) \\
  \code{venueAddress}  & \code{String}            & Dirección textual \\
  \code{location}      & \code{GeoJsonPoint}      & Coordenada lat/lon (índice 2dsphere) \\
  \idx{startsAt}       & \code{Instant}           & Fecha y hora de inicio \\
  \code{endsAt}        & \code{Instant}           & Fecha y hora de cierre \\
  \idx{genres}         & \code{List<String>}      & Géneros musicales \\
  \code{lineup}        & \code{List<LineupSlot>}  & Artistas embebidos (ver tabla siguiente) \\
  \code{price}         & \code{BigDecimal}        & Precio de entrada \\
  \code{capacity}      & \code{int}               & Capacidad máxima del venue \\
  \code{attendeeCount} & \code{int}               & Anotados (sincronizado desde Redis) \\
  \code{description}   & \code{String}            & Descripción del evento \\
  \idx{status}         & \code{String}            & \code{upcoming} / \code{live} / \code{finished} \\
  \code{city}          & \code{String}            & Ciudad \\
  \bottomrule
\end{tabularx}

{\footnotesize \idx{campo}\textsuperscript{*} = campo indexado}

\subsubsection{Subdocumento embebido: \code{LineupSlot}}

\begin{tabularx}{\textwidth}{@{}llX@{}}
  \toprule
  \textbf{Campo}  & \textbf{Tipo}   & \textbf{Descripción} \\
  \midrule
  \code{time}     & \code{String}   & Horario de inicio del slot (\emph{ej.} \code{"23:30"}) \\
  \code{slot}     & \code{String}   & Tipo de slot: \code{OPEN} / \code{PRIME} / \code{PEAK} / \code{CLOSE} \\
  \code{artist}   & \code{String}   & Artista o descripción del acto \\
  \bottomrule
\end{tabularx}

\vspace{0.5em}
\textbf{¿Por qué embebido y no referenciado?}
El lineup se lee siempre junto con el evento; nunca se consulta de forma
independiente. Embeber elimina un round-trip a la base de datos por cada
pantalla de detalle de evento.

\subsection{Índices y sus justificaciones}

\begin{tabularx}{\textwidth}{@{}lXl@{}}
  \toprule
  \textbf{Campo}  & \textbf{Justificación} & \textbf{Tipo} \\
  \midrule
  \code{name}        & Búsqueda por nombre (regex case-insensitive) & Regular \\
  \code{startsAt}    & Filtro temporal en el feed: \code{>= now()} & Regular \\
  \code{genres}      & Filtro por géneros en el feed & Multikey \\
  \code{status}      & Filtro \code{upcoming|live} en toda consulta de feed & Regular \\
  \code{location}    & Búsqueda \emph{near} por coordenadas & \code{2dsphere} \\
  \bottomrule
\end{tabularx}

\subsection{Desnormalización intencional}

\code{venueName} está repetido en cada evento aunque podría obtenerse
de una colección separada de venues. La decisión es intencional:
en la pantalla de feed, cada card del evento muestra el nombre del venue.
Si estuviera referenciado, necesitaríamos un \code{\$lookup} por cada
evento en el feed. La desnormalización convierte esa operación en lectura
directa del documento a costo de un pequeño riesgo de inconsistencia
(asumido como tolerable).

% ──────────────────────────────────────────────
\section{Redis — presencia en tiempo real}
\label{sec:redis}
% ──────────────────────────────────────────────

\subsection{Por qué Redis}

La presencia en tiempo real tiene requisitos que no puede cumplir ninguna
base de datos persistente:

\begin{itemize}[nosep]
  \item Latencia sub-milisegundo (un conteo en pantalla que sube en vivo)
  \item TTL nativo (la presencia expira sola sin cron jobs)
  \item Operaciones atómicas sobre sets y contadores
  \item Volumen de writes extremadamente alto durante picos nocturnos
\end{itemize}

Redis resuelve todo esto en memoria. El dato es efímero por naturaleza:
si la app se reinicia y pierde presencia, el estado se reconstruye en minutos
a medida que los usuarios re-interactúan.

\subsection{Esquema de claves}

Redis no tiene esquema formal. Las claves siguen convenciones de nombrado:

\begin{lstlisting}[style=redis, caption={patrones de clave en Redis}]
# en qué venue está un usuario ahora mismo
# tipo: String   TTL: 8 horas
presence:{userId}   →   "{venueId}"

# set de usuarios presentes en un venue
# tipo: Set   TTL: 8 horas
venue:{venueId}:present   →   { "userId1", "userId2", ... }

# contador de anotados para un evento
# tipo: String (entero)   TTL: ninguno (se borra manualmente)
event:{eventId}:attendees   →   "234"
\end{lstlisting}

\subsection{Operaciones del \code{PresenceService}}

\begin{tabularx}{\textwidth}{@{}lX@{}}
  \toprule
  \textbf{Método}             & \textbf{Operaciones Redis} \\
  \midrule
  \code{markPresent}          & \code{SET presence:\{u\} \{venueId\} EX 28800} + \code{SADD venue:\{v\}:present \{u\}} \\
  \code{markDeparted}         & \code{DEL presence:\{u\}} + \code{SREM venue:\{v\}:present \{u\}} \\
  \code{countPresent}         & \code{SCARD venue:\{v\}:present} \\
  \code{getPresentUsers}      & \code{SMEMBERS venue:\{v\}:present} \\
  \code{getCurrentVenue}      & \code{GET presence:\{u\}} \\
  \code{incrementAttendeeCount} & \code{INCR event:\{e\}:attendees} \\
  \bottomrule
\end{tabularx}

\vspace{0.5em}
\textbf{¿Por qué \code{SADD} para el set del venue?}
El set permite consultar el overlap entre \emph{``usuarios en venue X''} y
\emph{``amigos del usuario Y''} con una operación \code{SINTER} si en algún
momento se necesita esa query sin pasar por Neo4j.

\subsection{Sincronización con MongoDB}

\code{attendeeCount} en \code{EventDocument} refleja el contador de Redis.
La sincronización es eventual: Redis tiene el número en tiempo real; MongoDB
lo recibe periódicamente via \code{EventCatalogService.syncAttendeeCount()}.
El cliente siempre lee el conteo desde Redis si necesita precisión, y desde
MongoDB si tolera un lag de algunos minutos (ej. en listados estáticos).

% ──────────────────────────────────────────────
\section{Cassandra — historial y tendencias}
\label{sec:cassandra}
% ──────────────────────────────────────────────

\subsection{Por qué Cassandra}

El historial de check-ins es una serie temporal de escrituras continuas.
Las propiedades que necesita:

\begin{itemize}[nosep]
  \item Writes de alta frecuencia sin degradación de performance
  \item Lecturas siempre acotadas por un usuario o venue específico
  \item Sin necesidad de actualizar registros anteriores (append-only)
  \item Retención de datos por meses sin impacto en queries recientes
\end{itemize}

Cassandra fue diseñada exactamente para este patrón: optimiza para writes
sobre reads, y la clave de partición garantiza que cada query lea de una
sola partición (un solo nodo en el cluster).

\subsection{Tabla: \code{checkin\_history}}

\begin{lstlisting}[style=cypher, caption={CQL equivalente de la entidad}]
CREATE TABLE checkin_history (
    user_id     TEXT,
    checked_at  TIMESTAMP,
    event_id    TEXT,
    venue_id    TEXT,
    venue_name  TEXT,
    genre       TEXT,
    city        TEXT,
    PRIMARY KEY (user_id, checked_at)
) WITH CLUSTERING ORDER BY (checked_at DESC);
\end{lstlisting}

\begin{tabularx}{\textwidth}{@{}llX@{}}
  \toprule
  \textbf{Campo}     & \textbf{Rol en la clave}  & \textbf{Justificación} \\
  \midrule
  \pk{user\_id}      & Partition key             & Agrupa todo el historial de un usuario en una partición \\
  \ck{checked\_at}   & Clustering key (DESC)     & Ordena de más reciente a más antiguo; sin sort en query \\
  \code{event\_id}   & Columna                   & Referencia al evento en MongoDB \\
  \code{venue\_id}   & Columna                   & Para queries ``¿cuántas veces fui a este venue?'' \\
  \code{venue\_name} & Columna                   & Desnormalizado para display sin lookup \\
  \code{genre}       & Columna                   & Para análisis de afinidad musical histórica \\
  \code{city}        & Columna                   & Para análisis geográfico \\
  \bottomrule
\end{tabularx}

\vspace{0.5em}
\textbf{¿Por qué particionar por \code{user\_id}?}
La query más frecuente es \emph{``dame el historial reciente del usuario X''}.
Con \code{user\_id} como partition key, esa query toca exactamente una
partición. Si partiéramos por fecha o por venue, necesitaríamos leer
múltiples particiones para obtener el historial de un usuario.

\textbf{¿Por qué clustering descendente?}
El historial siempre se lee del más reciente al más antiguo. El orden DESC
en el clustering key hace que Cassandra almacene los datos ya ordenados
en ese sentido; no hay costo de reordenamiento en lectura.

\textbf{Nota sobre \code{ALLOW FILTERING}:}
La query \code{findByUserIdAndVenueId} usa \code{ALLOW FILTERING} porque
\code{venue\_id} no es parte de la primary key. Está admitido aquí porque
la partición ya está acotada por \code{user\_id} — el filtro aplica sobre
una cantidad pequeña de filas de una sola partición, no sobre toda la tabla.

\subsection{Tabla: \code{venue\_trends}}

\begin{lstlisting}[style=cypher, caption={CQL equivalente de VenueTrend}]
CREATE TABLE venue_trends (
    venue_id      TEXT,
    date          TEXT,
    hour          INT,
    checkin_count INT,
    unique_users  INT,
    top_genre     TEXT,
    PRIMARY KEY (venue_id, date, hour)
) WITH CLUSTERING ORDER BY (date DESC, hour ASC);
\end{lstlisting}

\begin{tabularx}{\textwidth}{@{}llX@{}}
  \toprule
  \textbf{Campo}        & \textbf{Rol}              & \textbf{Justificación} \\
  \midrule
  \pk{venue\_id}        & Partition key             & Un venue = una partición; toda su historia en un read \\
  \ck{date}             & Clustering (DESC)         & Días más recientes primero \\
  \ck{hour}             & Clustering (ASC)          & Dentro de un día, orden cronológico \\
  \code{checkin\_count} & Columna                   & Check-ins agregados en esa franja horaria \\
  \code{unique\_users}  & Columna                   & Usuarios únicos (aproximación) \\
  \code{top\_genre}     & Columna                   & Género más frecuente ese día/hora \\
  \bottomrule
\end{tabularx}

\vspace{0.5em}
Esta tabla responde preguntas como \emph{``¿a qué hora del viernes
Niceto tiene más gente?''} o \emph{``¿qué género convoca más en este venue?''}
sin escanear el historial completo de check-ins.

% ──────────────────────────────────────────────
\section{Flujo de check-in — orquestación de los cuatro motores}
\label{sec:checkin}
% ──────────────────────────────────────────────

El \code{CheckinService} es el único punto de la app que escribe en los
cuatro motores en una sola operación. El orden es por criticidad y latencia:

\begin{center}
\begin{tikzpicture}[
  box/.style={
    rectangle, rounded corners=3pt, minimum width=3.2cm, minimum height=0.75cm,
    font=\small\ttfamily, text centered
  },
  step/.style={
    rectangle, fill=codebg, text=white, minimum width=0.6cm, minimum height=0.6cm,
    font=\footnotesize\bfseries
  },
  arr/.style={-Stealth, thick, gray},
  note/.style={font=\footnotesize, text width=5cm}
]
  \node[box, fill=rediscolor!20, draw=rediscolor]
    (r) at (0,0) {Redis};
  \node[box, fill=neo4jcolor!20, draw=neo4jcolor]
    (n) at (0,-1.5) {Neo4j};
  \node[box, fill=casscolor!20, draw=casscolor]
    (c1) at (0,-3.0) {Cassandra (historial)};
  \node[box, fill=casscolor!20, draw=casscolor]
    (c2) at (0,-4.5) {Cassandra (tendencias)};

  \node[step] (s1) at (-2.2, 0)    {1};
  \node[step] (s2) at (-2.2,-1.5)  {2};
  \node[step] (s3) at (-2.2,-3.0)  {3};
  \node[step] (s4) at (-2.2,-4.5)  {4};

  \node[note, right=0.4cm of r]
    {Presencia inmediata.\\\code{markPresent} + \code{INCR attendees}};
  \node[note, right=0.4cm of n]
    {Relación \code{ATTENDING} para el motor de recomendación.};
  \node[note, right=0.4cm of c1]
    {Registro append-only en \code{checkin\_history}.};
  \node[note, right=0.4cm of c2]
    {Agrega en \code{venue\_trends} por día/hora.};

  \draw[arr] (s1) -- (r);
  \draw[arr] (s2) -- (n);
  \draw[arr] (s3) -- (c1);
  \draw[arr] (s4) -- (c2);

  \draw[arr] (r) -- (n);
  \draw[arr] (n) -- (c1);
  \draw[arr] (c1) -- (c2);
\end{tikzpicture}
\end{center}

\textbf{¿Por qué este orden?}
Redis primero porque el usuario ve la confirmación de presencia antes de
que el resto de los sistemas procesen. Neo4j segundo porque la relación
\code{ATTENDING} alimenta las recomendaciones de otros usuarios de forma
inmediata. Cassandra último porque el historial es auditoría — toleramos
un lag de milisegundos; no toleramos que un timeout de Cassandra bloquee
la confirmación al usuario.

\vspace{0.5em}
\textbf{Consideración de consistencia:}
Esta operación no es transaccional entre motores. Si Neo4j falla después de
que Redis ya escribió, el estado queda parcialmente consistente. Para el
caso de uso actual esto es aceptable: la presencia en Redis es la fuente
de verdad de corto plazo; Neo4j se puede reconstruir desde el historial de
Cassandra en un job batch si fuera necesario.

% ──────────────────────────────────────────────
\section{Resumen: qué lee cada pantalla}
\label{sec:pantallas}
% ──────────────────────────────────────────────

\begin{tabularx}{\textwidth}{@{}lX@{}}
  \toprule
  \textbf{Pantalla}       & \textbf{Fuente de datos} \\
  \midrule
  Feed principal
    & MongoDB: eventos activos/próximos en la ciudad. Redis: contadores de
      anotados en tiempo real. Neo4j: qué amigos van a cada evento (overlay). \\
  \addlinespace
  Detalle de evento
    & MongoDB: toda la info del evento (lineup, descripción, venue). Redis:
      conteo exacto de asistentes. Neo4j: amigos que se anotaron y su estado
      (\code{confirmed/pending}). \\
  \addlinespace
  Mapa de jodas activas
    & MongoDB: coordenadas y nombre de eventos con \code{status=live}.
      Redis: conteos por venue. \\
  \addlinespace
  Pantalla de pings / notificaciones
    & Redis pub/sub (eventos de presencia en tiempo real).
      Neo4j: validación de amistad para filtrar notificaciones. \\
  \addlinespace
  Historial personal
    & Cassandra: \code{checkin\_history} particionado por \code{user\_id}. \\
  \addlinespace
  Recomendaciones
    & Neo4j: query de 2 saltos (amigos → eventos). MongoDB: enriquecimiento
      de los \code{eventId} devueltos por el grafo. \\
  \bottomrule
\end{tabularx}

% ──────────────────────────────────────────────
\section{Puertos y configuración}
\label{sec:config}
% ──────────────────────────────────────────────

\begin{tabularx}{\textwidth}{@{}llll@{}}
  \toprule
  \textbf{Motor}     & \textbf{Puerto app} & \textbf{Puerto admin UI} & \textbf{Imagen Docker} \\
  \midrule
  Neo4j              & \code{7687} (Bolt) / \code{7474} (HTTP) & \code{7474} (browser) & \code{neo4j:5} \\
  MongoDB            & \code{27017}        & \code{8081} (Mongo Express) & \code{mongo:7} \\
  Redis              & \code{6379}         & \code{5540} (RedisInsight)  & \code{redis:7-alpine} \\
  Cassandra          & \code{9042}         & —                           & \code{cassandra:4.1} \\
  \bottomrule
\end{tabularx}

\vspace{0.5em}
Las herramientas de administración visual (Mongo Express, RedisInsight)
se levantan por separado con:

\begin{lstlisting}[style=redis]
docker compose --profile admin up -d
\end{lstlisting}

% ──────────────────────────────────────────────
\section*{Apéndice: decisiones descartadas}
% ──────────────────────────────────────────────

\begin{description}[leftmargin=2cm, style=nextline]
  \item[PostgreSQL para todo]
    Las queries de recomendación requerirían JOINs recursivos (CTEs)
    que no escalan bien con el crecimiento de la red social. El índice
    geográfico de PostGIS funciona pero agrega complejidad operacional.

  \item[Un solo MongoDB para todo]
    MongoDB puede hacer grafos con \code{\$graphLookup} pero con un límite
    de 100MB de RAM y sin el nivel de optimización de un motor dedicado.
    Los writes de alta frecuencia del historial tampoco son el punto fuerte
    de MongoDB.

  \item[DynamoDB en lugar de Cassandra]
    DynamoDB resolvería el mismo caso pero es managed cloud (vendor lock-in).
    Cassandra es self-hosted y portable entre ambientes. Para un proyecto
    académico/portfolio, Cassandra aporta más valor de aprendizaje.

  \item[Elasticsearch para el feed]
    La búsqueda full-text de eventos no es el core del producto en esta
    fase. Los índices de MongoDB son suficientes. Si en el futuro se
    agregan búsquedas complejas de texto o faceting, Elasticsearch sería
    la extensión natural.
\end{description}

\end{document}
